% from ega4-2-fr.tex
\oldpage[IV-2]{102}

les corps résiduels correspondants sont $k(\mathfrak m)=k$ et
$k(\mathfrak m')=k(S)$, corps des fractions de $V$; en vertu du choix de
$k$, ces corps sont isomorphes. Désignons par $\varepsilon$ et
$\varepsilon'$ les homomorphismes canoniques de $E$ sur
$E/\mathfrak m=k(\mathfrak m)$ et $E/\mathfrak m'=k(\mathfrak m')$
respectivement; soit d'autre part $\sigma$ un isomorphisme de
$k(\mathfrak m)$ sur $k(\mathfrak m')$, et considérons le sous-anneau $C$
de $E$ formé des $x\in E$ tels que
$\varepsilon'(x)=\sigma(\varepsilon(x))$ (cette construction est un cas
particulier des « procédés de recollement » qui seront étudiés
systématiquement au chap. V). Il est immédiat que
$\mathfrak n=\mathfrak m\mathfrak m'=\mathfrak m\cap\mathfrak m'$ est un
idéal maximal de $C$, $C/\mathfrak m\mathfrak m'$ s'identifiant au
sous-corps de $(E/\mathfrak m)\times(E/\mathfrak m')$ formé des couples
$(z,\sigma(z))$. On a $E=C+C(ST-1)$, autrement dit $E$ est une
$C$-algèbre finie et est évidemment la clôture intégrale de $C$; nous allons
voir en outre que $C$ est noethérien: cela résultera du lemme suivant:

\begin{lemma}[5.6.11.1]
\EGAReaderTarget{ega4.c04.lem.5-6-11-1}{IV.2.s5.6.11.1-fr}
Soient $R$ un anneau, $S$ un sous-anneau,
$\mathfrak K=\operatorname{Ann}_S(R/S)$ le conducteur de $S$ dans $R$
(plus grand idéal de $S$ qui est aussi un idéal de $R$).
\begin{enumerate}
\item[(i)] Pour tout idéal $\mathfrak J\subset\mathfrak K$ de $R$, il existe
une application strictement croissante de l'ensemble des idéaux
$\mathfrak a$ de $S$ tels que $R\mathfrak a=\mathfrak J$ sur l'ensemble des
sous-$(S/\mathfrak K)$-modules de
$\mathfrak J/\mathfrak K\mathfrak J$.
\item[(ii)] Si $S/\mathfrak K$ et $R$ sont noethériens et si $R$ est un
$S$-module de type fini, alors toute suite croissante d'idéaux de $S$
contenus dans $K$ est stationnaire.
\end{enumerate}
\end{lemma}

(i) On a en effet,
$\mathfrak K\mathfrak J=\mathfrak K(R\mathfrak a)
=\mathfrak K\mathfrak a\subset\mathfrak a$
($\mathfrak a$ étant un idéal de $S$), d'où la conclusion.

(ii) Soit $(\mathfrak L_n)$ une suite croissante d'idéaux de $S$ contenus
dans $\mathfrak K$. La suite des idéaux $R\mathfrak L_n$ de $R$ est
stationnaire puisque $R$ est noethérien; on peut donc supposer que tous les
idéaux $R\mathfrak L_n$ sont égaux à un même idéal $\mathfrak J$ de $R$.
Comme $R$ est noethérien, $\mathfrak J/\mathfrak K\mathfrak J$ est un
$R$-module de type fini, donc aussi un $S$-module de type fini, et par suite
un $(S/\mathfrak K)$-module de type fini. Mais comme $S/\mathfrak K$ est
noethérien, $\mathfrak J/\mathfrak K\mathfrak J$ est un
$(S/\mathfrak K)$-module noethérien, et la conclusion résulte de (i).

Nous allons appliquer ce lemme en prenant $R=E$, $S=C$; il est clair que
$\mathfrak n$ est le conducteur de $C$ dans $E$; en outre, pour tout idéal
$\mathfrak a$ de $C$,
$\mathfrak a/(\mathfrak n\cap\mathfrak a)$ est isomorphe à
$(\mathfrak a+\mathfrak n)/\mathfrak n$, donc un sous-$C$-module de
$C/\mathfrak n$, qui est un $C$-module simple; il suffit donc de démontrer
que tout idéal $\mathfrak a\cap\mathfrak n$ de $C$ est de type fini, et cela
résulte de (5.6.11.1).

Il résulte de (0, 16.1.5) que l'on a $\dim(C)=\dim(E)=2$. Prenons
$A=C_{\mathfrak n}$; si $U=C-\mathfrak n$, l'anneau de fractions
$B=U^{-1}E$ est donc la clôture intégrale de $A$, et c'est un $A$-module de
type fini; d'ailleurs, comme $\mathfrak m$ et $\mathfrak m'$ sont les seuls
idéaux premiers de $E$ contenant $\mathfrak n$, $B$ est un anneau semi-local
dont les anneaux locaux en les deux idéaux maximaux sont isomorphes à
$E_{\mathfrak m}$ et $E_{\mathfrak m'}$ respectivement, donc sont
respectivement de dimension 2 et 1. Cela montre que $\dim(B)=2$, donc aussi
$\dim(A)=2$ (0, 16.1.5). Comme $A$ est un anneau local intègre, il est
nécessairement caténaire (tout idéal premier distinct de 0 et de l'idéal
maximal étant nécessairement de hauteur 1); mais il ne vérifie pas la
conclusion de (5.6.10), et \emph{a fortiori} n'est pas universellement
caténaire.

Notons encore que $\mathfrak n$ est un idéal de hauteur 2 dans $C$, et que
pour tout idéal premier $\mathfrak p$ de hauteur 1 dans $C$, il existe un
seul idéal premier $\mathfrak p'$ de $E$ au-dessus de $\mathfrak p$,
nécessairement de hauteur 1, et tel que
$C_{\mathfrak p}=E_{\mathfrak p'}$. En effet, il y a au moins un idéal
premier $\mathfrak p'$ de $E$ au-dessus de $\mathfrak p$ et il résulte du
théorème de Cohen-Seidenberg qu'un

